\subsection{Local-code outer interface}
\label{sec:local-code-outer}

The outer side replaces sliding random bands by short local codes supplied
through explicit spectra or certified spectrum envelopes.  The local code is an
input to the theorem; RM, BCH, or another code with a certified spectrum are
possible choices.

\subsubsection{Direct-sum outer spectrum}
\label{subsec:local-code-outer-interface}

The structured outer is a direct sum of short local codes.  The proof interface
is deliberately local-code independent: the outer contributes a weight
enumerator, the interleaver turns each outer word of weight \(h\) into a uniform
slice input, and the inner contributes an upper bound
\(p_h^{\rm in}(\delta)\) for that slice.  Systematic form is not required for
the outer code.

\begin{lemma}[Block-outer first-moment interface]
\label{lem:block-outer-first-moment-interface}
Let \(C_{\mathrm{loc}}\subseteq\F_2^{n_0}\) be a binary linear \([n_0,b]\) code
with homogeneous weight enumerator
\[
W_{\mathrm{loc}}(z)=\sum_j A^{\mathrm{loc}}_j z^j,
\]
and let \(C_{\mathrm{out}}=C_{\mathrm{loc}}^{\oplus M}\).  Its message length
is \(k=Mb\), its outer/interleaver length is \(N=Mn_0\), and its weight
enumerator is
\[
W_{\mathrm{out}}(z)=W_{\mathrm{loc}}(z)^M .
\]
Write \(A_h^{\mathrm{out}}=[z^h]W_{\mathrm{out}}(z)\), omitting the all-zero
word from the nonzero-codeword sum.  If \(p_h^{\rm in}(\delta)\) is any valid
upper bound on the chosen inner construction producing output weight at most
\(\lfloor\delta N\rfloor\) from an interleaved outer word of weight \(h\), then
\[
\Pr\Bigl[
 \Ctwo\circ\Pi\big|_{C_{\mathrm{out}}}\text{ is not injective}
 \ \text{or}\ d_{\min}(\Csc)\le \lfloor\delta N\rfloor
\Bigr]
\le
\mu_{\mathrm{block}}(\delta)
:=
\sum_{h=1}^{N} A_h^{\mathrm{out}}p_h^{\rm in}(\delta).
\]
\end{lemma}

\begin{proof}
The direct-sum enumerator identity follows by multiplying the local generating
functions block by block.  Conditional on any nonzero outer word of weight
\(h\), the uniform interleaver randomizes only the support positions seen by
the inner code, so the same inner bound \(p_h^{\rm in}(\delta)\) applies to
every such word.  Summing these failure probabilities over all nonzero outer
words gives \(\mu_{\mathrm{block}}(\delta)\).  Markov's inequality bounds the
event that the restriction has a nonzero kernel word or a nonzero image word
below the target distance.
\end{proof}

\begin{lemma}[Subcode spectrum domination]
\label{lem:block-outer-subcode-domination}
Let \(C'\subseteq C_{\mathrm{out}}\) be a linear subcode, with weight
enumerators \(A'_h\) and \(A_h^{\mathrm{out}}\), respectively.  Then
\(A'_h\le A_h^{\mathrm{out}}\) for every \(h\).  If in addition
\(d_{\min}(C')\ge d_0\), then the first-moment bound for \(C'\) is at most
\[
 \sum_{h=d_0}^{N}A_h^{\mathrm{out}}p_h^{\rm in}(\delta).
\]
\end{lemma}

\begin{proof}
Every weight-\(h\) word of \(C'\) is a distinct weight-\(h\) word of
\(C_{\mathrm{out}}\), proving coefficientwise domination.  The distance
hypothesis makes \(A'_h=0\) below \(d_0\); applying
Lemma~\ref{lem:block-outer-first-moment-interface} and then replacing each
remaining \(A'_h\) by \(A_h^{\mathrm{out}}\) gives the claim.
\end{proof}

\begin{lemma}[Distance under scalar extension]
\label{lem:binary-code-scalar-extension}
Let \(G\in\F_2^{K\times N}\) generate a binary linear \([N,K,d]\) code.
For every \(s\ge1\), interpreting the same matrix over \(\F_{2^s}\) generates
an \(\F_{2^s}\)-linear \([N,K,d]\) code in symbol Hamming distance.
\end{lemma}

\begin{proof}
The binary row rank of \(G\) remains \(K\) over every extension field.  Fix an
\(\F_2\)-basis \(\beta_1,\ldots,\beta_s\) of \(\F_{2^s}\).  Every extended-field
message has a unique decomposition \(x=\sum_j\beta_jx_j\), with
\(x_j\in\F_2^K\), and hence
\[
 xG=\sum_j\beta_j(x_jG).
\]
If \(xG\ne0\), some binary component \(x_jG\) is nonzero.  Its support is
contained in the symbol support of \(xG\), so \(\wt(xG)\ge d\).  Conversely,
a binary codeword of weight \(d\) is also an extended-field codeword, proving
equality of the two minimum distances.
\end{proof}

\subsubsection{Finite certificate template}
\label{subsec:local-code-certificate-template}

\begin{lemma}[Finite block-certificate template]
\label{lem:finite-block-certificate-template}
In the setting of Lemma~\ref{lem:block-outer-first-moment-interface}, suppose
the nonzero outer weights are partitioned into a finite list of audited row
families \(\mathcal R\).  For each row \(R\in\mathcal R\), suppose a verifier or
proof establishes
\[
\sum_{h\in R} A_h^{\mathrm{out}}p_h^{\rm in}(\delta)\le 2^{-\lambda_R}.
\]
If the rows cover every nonzero outer word and every inner placement/episode
case used by the bound, then
\[
\Pr\Bigl[
 \Ctwo\circ\Pi\big|_{C_{\mathrm{out}}}\text{ is not injective}
 \ \text{or}\ d_{\min}(\Csc)\le\lfloor\delta N\rfloor
\Bigr]
\le
\sum_{R\in\mathcal R}2^{-\lambda_R}.
\]
In particular, proving the right-hand side is at most \(2^{-\lambda}\) gives
failure probability at most \(2^{-\lambda}\).
\end{lemma}

\begin{proof}
This is bookkeeping after Lemma~\ref{lem:block-outer-first-moment-interface}.
The row hypotheses form a union-bound cover of the first-moment sum over all
nonzero outer words and all inner cases, and summing the row bounds gives the
claim.
\end{proof}
